\documentclass[xcolor=pdftex, x11names,table, oneside,11pt]{book}
%\usepackage{comment}
%\usepackage[utf8x]{inputenc}
%\usepackage[spanish, es-lcroman,es-tabla]{babel}
%\usepackage{amsmath}
%\usepackage{amsfonts}
%\usepackage{amssymb}
%\usepackage{graphicx}
\usepackage[letterpaper]{geometry}
\geometry{verbose,tmargin=2.5cm,bmargin=2.5cm,lmargin=3.5cm,rmargin=2.5cm}
%\setcounter{tocdepth}{3}
\usepackage{textcomp}
%\usepackage{amstext}
%\usepackage{setspace}
\usepackage[toc, page]{appendix}
\usepackage{tikz}
%\usetikzlibrary{arrows.meta}
%\usepackage{pgf}
%\usepackage{pgfplots}
%\usepgflibrary{decorations.pathmorphing}
\usepackage{xcolor}
\usepackage{tabulary}
\usepackage{colortbl}
%\usetikzlibrary{patterns}
%\usepackage{tikz-3dplot}
%\usepackage{etoolbox}
\usepackage{hyperref}
\usepackage{rotating}
\usepackage{float}
%\usepackage[section]{placeins}
\usepackage{afterpage}
%\usepackage{cleveref}
\usepackage{array}
%\usepackage{titlesec}
\usepackage{pdfpages}
\usepackage{hyperref}
%\usepackage[Sonny]{fncychap}
%\usepackage{caption}
%\usepackage[utopia]{quotchap}
%\hyphenation{de-sa-pa-ri-ci\'on}
%\onehalfspace
\usepackage{titletoc}
%\usepackage{titlesec}
\usepackage[shortlabels]{enumitem}
\usepackage{lipsum}
\hyphenation{trabajos}
%\titleformat*{\section*}{\color{c0}\normalfont\bfseries\LARGE}
\usepackage{booktabs}
\usepackage{multirow} 
\usepackage{multicol} 
\usepackage{adjustbox}
\usepackage[most]{tcolorbox}
\renewcommand{\arraystretch}{1.5}

%\usepackage{roboto}
%\usepackage{draftwatermark}
%\SetWatermarkText{\textbf{\textsc{EN REVISION}}}
%\SetWatermarkLightness{0.7}
%\SetWatermarkScale{2.5}
%\DraftwatermarkOptions{color={[rgb]{0.98, 0.91, 0.71}}}

%\newcommand{\art}[1]{{\section{Artículo \thesection  #1}}
\newcounter{articulo}
\refstepcounter{articulo}

\newcounter{capitulo}
\refstepcounter{capitulo}

\newcommand{\art}[1]{
\addcontentsline{toc}{section}{Art\'iculo \thearticulo. #1} 
\section*{Artículo \thearticulo: #1} 
\refstepcounter{articulo}
}

%%\addcontentsline{toc}{section}{Art\'iculo \thearticulo. Alcance}


\newcommand{\Chapter}[1]{
\addcontentsline{toc}{chapter}{Cap\'itulo \thecapitulo: #1} 
\chapter*{Cap\'itulo \thecapitulo: #1}
\refstepcounter{capitulo}
}

\definecolor{c1}{RGB}{143,85,69}
\definecolor{c2}{RGB}{218,37,39}
\definecolor{c3}{RGB}{231,121,25}
\definecolor{c4}{RGB}{247,215,75}
\definecolor{c5}{RGB}{0,146,64}
\definecolor{c6}{RGB}{0,141,221}
\definecolor{c7}{RGB}{153,71,121}
\definecolor{c8}{RGB}{171,169,169}
\definecolor{c9}{RGB}{255,255,255}
\definecolor{c0}{RGB}{0,0,0}
\definecolor{r1}{RGB}{246,205,162}
\definecolor{azul1}{RGB}{15,70,124}
\definecolor{azul2}{RGB}{63,108,150}
\definecolor{r2}{RGB}{232,54,67}
% \dottedcontents{<section>}[<left>]{<above-code>} {<label width>}{<leader width>}
\dottedcontents{chapter}[0em]{\addvspace{1pc} \bfseries}{1em}{1pc}%[\addvspace{1pc}]
\dottedcontents{section}[2em]{}{2em}{1pc}
\dottedcontents{subsection}[5em]{}{3em}{1pc}
\renewcommand{\contentsname}{\centering Contenidos}

\makeatletter

\usepackage{fancyhdr}
\pagestyle{fancy}
\pagestyle{myheadings}

\renewcommand{\appendixname}{Apéndices}
\renewcommand{\chaptername}{Cap\'itulo}

\renewcommand{\appendixtocname}{Apéndices}
\renewcommand{\appendixpagename}{Apéndices}
\newcommand{\itcr}{Instituto Tecnol\'ogico de Costa Rica}
\newcommand{\tfg}{Trabajo Final de Graduaci\'on}
\newcommand{\tfgs}{Trabajos Finales de Graduaci\'on}
\newcommand{\ef}{Escuela de F\'isica}
\newcommand{\ualif}{Unidad Acad\'emica de la Licenciatura en Ingenier\'ia F\'isica}
\newcommand{\cualif}{Consejo de \ualif}
\newcommand{\lif}{Licenciatura en Ingenier\'ia F\'isica}
\newcommand{\Def}[2]{\textbf{\texttt{#1:}}#2 \\[.25cm]}
%\newcommand{\Def}[2]{\begin{minipage}{\textwidth}\\ \texttt{#1:}#2  \end{minipage} }

\makeatletter
\renewenvironment{titlepage}
 {%
  \if@twocolumn
    \@restonecoltrue\onecolumn
  \else
    \@restonecolfalse\newpage
  \fi
  \thispagestyle{empty}%
 }
 {%
  \if@restonecol
    \twocolumn
  \else
    \newpage
  \fi
 }
 
%\titleformat{\chapter}[display]{\normalfont\huge\bfseries}{\chaptertitlename\ \thechapter}{0pt}{\Huge}

%\titleformat{\section}[box]{\Huge\bfseries}{-}{0pt}{\Huge\bfseries}

%\titlespacing*{\chapter}{0pt}{-1.5cm}{1cm}


\setcounter{secnumdepth}{3}
\setcounter{tocdepth}{3}

\hypersetup{
    colorlinks=true,
    linkcolor=blue,
    filecolor=magenta,      
    urlcolor=cyan,
    pdfpagemode=FullScreen,
    }

%glosario
%\makeglossaries
\setlength{\parindent}{0pt}

\makeatletter
\renewcommand{\@seccntformat}[1]{\llap{\textcolor{c1}{\csname the#1\endcsname}\hspace{1em}}}                    
\renewcommand{\section}{\@startsection{section}
{1}{0.9cm}%0.9 corre la numeraci�n hacia adentro
{-4ex \@plus -1ex \@minus -.4ex}
{1ex \@plus.2ex }
{\color{c6}\normalfont\Large\sffamily\bfseries}}
\renewcommand{\subsection}{\@startsection {subsection}{2}{\z@}
{-3ex \@plus -0.1ex \@minus -.4ex}
{0.5ex \@plus.2ex }
{\color{c2}\normalfont\large\sffamily\bfseries}}
\renewcommand{\subsubsection}{\@startsection {subsubsection}{3}{\z@}
{-2ex \@plus -0.1ex \@minus -.2ex}
{0.2ex \@plus.2ex }
{\color{c3}\normalfont\small\sffamily\bfseries}}                        
\renewcommand\paragraph{\@startsection{paragraph}{4}{\z@}
{-2ex \@plus-.2ex \@minus .2ex}
{0.1ex}
{\normalfont\small\sffamily\bfseries}}
\makeatother
% Fin numeraci�n secciones



%---------------------------------------------------------------------------------
%  Entornos:  Ejemplo, teorema, proposici�n, lema, lista de ejercicios, 
%             caja interludio, caja simple  
%---------------------------------------------------------------------------------

%  Cajas con el paquete  tcbcolor
%  CONTADORES: ejemplo, definicion, lema, teorema, corolario, proposicion,ejercicio 
\newcounter{tcbteo}[chapter]
\renewcommand{\thetcbteo}{\thechapter.\arabic{tcbteo}}

\newcounter{tcbdefi}[chapter]
\renewcommand{\thetcbdefi}{\thechapter.\arabic{tcbdefi}}

\newcounter{tcblema}[chapter]
\renewcommand{\thetcblema}{\thechapter.\arabic{tcblema}}

\newcounter{tcbcoro}[chapter]
\renewcommand{\thetcbcoro}{\thechapter.\arabic{tcbcoro}}

%  \newcounter{tcbListaEjercicios}[chapter]
%  \renewcommand{\thetcbListaEjercicios}{\thechapter.\arabic{tcbListaEjercicios}}

\newcounter{tcbpropo}[chapter]
\renewcommand{\thetcbpropo}{\thechapter.\arabic{tcbpropo}}

\newcounter{tcbvoca}[chapter]
\renewcommand{\thetcbvoca}{\thechapter.\arabic{tcbvoca}}

\newlength{\examlen}
\tikzset{
    wnodeTeorema/.style={%
         rectangle,  top color=gray!5, bottom color=gray!5,
         inner sep=1mm,anchor=west,font=\small\bf\sffamily},
   wnodeminimo/.style={%
         rectangle,  top color=white, bottom color=white,
         text=azulF,inner sep=1mm,anchor=west,font=\small\bf\sffamily}      
}

%\begin{teorema}  o \begin{teorema}[de tal] o \begin{teorema}[][ref]
% Teorema -----------------------------------------------------
\newtcolorbox{wwteorema}[2][]{%
arc=0mm,breakable,enhanced,colback=gray!5,boxrule=0pt,top=7mm,drop fuzzy shadow,
fontupper=\normalsize,step and label={tcbteo}{pink},
overlay unbroken = {\draw[color=pink,line width=0.2pt] (frame.north west)--([xshift=0pt]frame.north east);
%Caja de T�tulo: teo --
\node[wnodeTeorema](tituloteo) at ([xshift=0pt, yshift=-4mm]frame.north west)
{\textbf{\color{pink} #2}};
%Borde superior --
\draw[pink,line width=2.5cm] ([xshift=1.25cm, yshift=0cm]frame.north west)-- +(\examlen,3pt);
},%
overlay first = {\draw[color=pink,line width=0.2pt] (frame.north west)--([xshift=0pt]frame.north east);
%Caja de T�tulo: teo --
\node[wnodeTeorema](tituloteo) at ([xshift=0pt, yshift=-4mm]frame.north west)
{\textbf{\color{pink} #2}};
%Borde superior --
\draw[pink,line width=2.5cm] ([xshift=1.25cm, yshift=0cm]frame.north west)-- +(\examlen,3pt);
},%
% Mantener borde en cambio de p�gina 
overlay last = {\draw[color=pink,line width=0.2pt] (frame.south west)--([xshift=0pt]frame.south east);},
#1}
%-

